% ch7_a 第7章 §7.1 (书页 275–287 / p290–p302)
% 第7章 THEORY OF QUANTUM HALL STATES
% 块界说明：本块自章首起，无 A-TAIL/JOIN。
% 书页 287（p302）末句 “...then it is easy to” 跨页延续至书页 288（p303）顶部
% “understand why it has a filling fraction ν=1/3.”，按约定半句译文留在本文件正文最后，
% 由下一块（ch7_b）的 A-TAIL 接续。
\chapter{量子霍尔态理论}

\begin{keypoints}
\item 分数量子霍尔液体代表着一些新的物质态，它们含有极为丰富的内部结构。
\item 不同的分数量子霍尔液体无法用序参量和局域算符的长程序来刻画。我们需要一套全新的理论来描述分数量子霍尔液体。
\item 分数量子霍尔液体含有无能隙的边缘激发，这些边缘激发构成所谓的手征 Luttinger 液体。边缘处的手征 Luttinger 液体也具有非常丰富的结构，这是体内丰富内部结构的一种反映。
\end{keypoints}

分数量子霍尔（FQH）效应出现在强磁场中的二维电子系统里。自 1982 年被发现以来（Tsui et al., 1982; Laughlin, 1983），对 FQH 系统的实验不断揭示出许多新现象和意外之处。所观察到的丰富分级（hierarchical）结构（Haldane, 1983; Halperin, 1984）表明，展示分数量子霍尔效应的电子态（这些态称为 FQH 液体）含有极为丰富的内部结构。事实上，FQH 液体代表着一种全新的物质态，它与前几章讨论过的对称性破缺态非常不同。要理解这种新的态，人们需要发展新的概念和新的技术。在本章中，我们将介绍 FQH 态的一般理论。对 FQH 态中新序（即拓扑序）的详细讨论将在第 8 章给出。

\section{阿哈罗诺夫--玻姆效应与分数统计}

\subsection{阿哈罗诺夫--玻姆效应——不接触而使粒子偏转}

\begin{keypoints}
\item 可缩回路（contractible loop）的局域相位与不可缩回路（non-contractible loop）的整体相位。
\item 非零的局域相位代表一种力，而非零的整体相位不对应任何经典的力。
\end{keypoints}

%FIG{7.1}: {挖去了 $\pmb{r}=0$ 点的平面，在 $\pmb{r}=0$ 处有一根磁通管。回路 $A$ 不可缩，缠绕数 $n_w=1$；回路 $B$ 可缩。}

为了给以后关于量子霍尔（QH）态中量子干涉的讨论做准备，让我们先简短地讨论一下阿哈罗诺夫--玻姆（A--B）效应。首先，我们想引入两个概念，即局域相位（local phase）与整体相位（global phase）。局域相位与可缩回路相联系。把一个带电粒子沿一条可缩回路缓慢地移动一周，会产生如下的局域相位：
\begin{equation*}
\text{局域相位} = \exp\left(\mathrm{i}e\oint\pmb{A}\,\mathrm{d}\pmb{x}\right) = \exp\left(\mathrm{i}e\int\pmb{B}\cdot\mathrm{d}\pmb{S}\right).
\end{equation*}
磁场 $\pmb{B}=0$ 当且仅当所有可缩回路的局域相位都为零。因此，不受磁力作用的带电粒子没有局域相位。

整体相位则与不可缩回路相关。为了在一个简单的背景中理解整体相位，让我们考虑一个在二维平面上运动的带电粒子。我们挖去 $\pmb{r}=0$ 点，这改变了空间的拓扑。然后在 $\pmb{r}=0$ 处放置一根无穷细的磁通管，其磁通为 $\Phi$（见图 7.1）。当粒子绕 $\pmb{r}=0$ 一周时，会产生一个相位 $\mathrm{e}^{\mathrm{i}e\Phi}$。这样的相位称为整体相位。更一般地，
\begin{gather*}
\text{整体相位} = \mathrm{e}^{\mathrm{i}e\Phi n_w}\\
n_w = \text{缠绕数}.
\end{gather*}

显然，尽管整体相位非零，粒子却感受不到任何磁力，因为在远离 $\pmb{r}=0$ 处没有磁场。我们看到，一个自由粒子（即不受力的粒子）可以具有非零的整体相位。

只有当不可缩回路存在时（即空间非单连通时），整体相位才会存在。非单连通空间的另一个例子是环面（见图 7.2）。环面上的自由粒子同样可以具有整体相位。

上述自由粒子（在 $\pmb{r}=0$ 处有一根磁通管）可以用如下哈密顿量描述：
\begin{equation*}
H = -\frac{1}{2m}\left(\pmb{\partial}-\mathrm{i}\pmb{a}\right)^2, \qquad (a_x, a_y) = \frac{1}{r^2}(-y, -x)\frac{\phi}{2\pi}
\end{equation*}

%FIG{7.2}: {带有两个不可缩回路 $A$ 和 $B$ 的环面。}

%FIG{7.3}: {一束粒子穿过一排管子，每根管子中都有磁通。}

场强为零，即
\begin{equation*}
\pmb{b} = \partial_x a_y - \partial_y a_x = 0 \ \Longrightarrow\ \text{局域相位} = 0
\end{equation*}
对于绕 $\pmb{x}=0$ 一周的回路，我们有
\begin{equation*}
\oint\pmb{a}\cdot\mathrm{d}\pmb{x} = \phi \ \Longrightarrow\ \text{整体相位} \neq 0\ .
\end{equation*}

非零的局域相位会影响经典的运动方程。自旋的贝里相就是局域相位的一个例子（见 2.3 节）。与之相反，整体相位不会影响经典的运动方程，但它会影响粒子的量子性质。

\subsection*{问题 7.1.1}

\textbf{不接触而使其偏转}\quad 尽管整体相位不产生任何经典的力，它仍然能使一个运动的粒子偏转。考虑图 7.3 中的装置：一束电荷为 $e$、动量为 $\pmb{k}$ 的粒子穿过一排不可穿透的管子。若每根管子中都有磁通 $\Phi$，试计算偏转角 $\theta$。

%FIG{7.4}: {$V_-$ 空间：挖去 $\pmb{r}=0$ 点，并把 $\pmb{r}$ 与 $-\pmb{r}$ 认同为同一个点。}

\subsection{具有硬芯条件的粒子与分数统计}

\begin{keypoints}
\item 具有硬芯条件的粒子，其位形空间不是单连通的。位形空间的非平凡拓扑允许非平凡的统计。
\item 分数统计存在，而且只存在于二维空间之中。
\end{keypoints}

A--B 效应的一个具体体现是二维中的分数统计（Leinaas and Myrheim, 1977; Wilczek, 1982）。考虑二维中两个自由的、具有硬芯的全同粒子。\footnote{这里，所谓自由粒子，是指彼此分开时不受任何力作用的粒子；但这些粒子可以受到硬芯条件的约束。}其位形空间由下式给出：
\begin{equation*}
\begin{aligned}
\text{位形空间} &= \left\{(\pmb{r}_1, \pmb{r}_2)\big|_{\pmb{r}_1\neq\pmb{r}_2,\ (\pmb{r}_1,\pmb{r}_2)\sim(\pmb{r}_2,\pmb{r}_1)}\right\}\\
&= \left\{(\pmb{r}_+, \pmb{r}_-)\big|_{\pmb{r}_-\neq 0,\ (\pmb{r}_+,\pmb{r}_-)\sim(\pmb{r}_-,-\pmb{r}_-)}\right\} = V_+\otimes V_-.
\end{aligned}
\end{equation*}
其中 $\pmb{r}_+=\frac{1}{2}(\pmb{r}_1+\pmb{r}_2)$，$\pmb{r}_-=\pmb{r}_1-\pmb{r}_2$，$V_+=\{\pmb{r}_+\}$，而 $V_-=\{\pmb{r}_-|_{\pmb{r}_-\neq 0,\ \pmb{r}_-\sim-\pmb{r}_-}\}$（见图 7.4）。这里 $V_+$ 是通常的二维平面，而 $V_-$ 只是二维平面的一半，因为 $\pmb{r}_-$ 与 $-\pmb{r}_-$ 被视为同一个点。同时，$\pmb{r}_-=0$ 点已从 $V_-$ 中挖去。

如上一节所指出的，“自由”意味着局域相位为零。由于两粒子位形空间不是单连通的（$V_-$ 非单连通），整体相位存在，并且我们有权自由选择整体相位。注意，$V_-$ 中的不可缩回路由缠绕数（winding number）$n_w$ 标记（图 7.5）。于是，我们可以给不可缩回路指派如下整体相位：
\begin{equation*}
\text{整体相位} = \mathrm{e}^{\mathrm{i}\theta n_w}
\end{equation*}

现在很清楚，在量子力学中存在不同种类的自由全同粒子，由一个参数 $\theta$ 标记。

%FIG{7.5}: {$V_-$ 空间中缠绕数 $n_w=1$ 与 $n_w=2$ 的两个回路。}

\begin{itemize}
\item 参数 $\theta$ 描述多粒子位形空间中的磁通（即整体相位）。
\item 参数 $\theta$ 决定统计性质。
\end{itemize}

注意，$n_w=1$ 的回路把 $\pmb{r}_-$ 连接到 $-\pmb{r}_-$，即把 $(\pmb{r}_1,\pmb{r}_2)$ 连接到 $(\pmb{r}_2,\pmb{r}_1)$，这对应于交换两个粒子。于是，粒子的统计如下：
\begin{equation*}
\left\{\begin{array}{ccc}
\theta = 0 & \rightarrow & \text{玻色子}\\
\theta = \pi & \rightarrow & \text{费米子}\\
\theta = \text{其他} & \rightarrow & \text{任意子}
\end{array}\right.
\end{equation*}

上述两个自由全同粒子（即两个具有统计 $\theta$ 的任意子（anyon））的哈密顿量为
\begin{gather*}
H = -\frac{1}{2m}(\partial_{\pmb{r}_+})^2 - \frac{1}{2}\frac{1}{m}(\partial_{\pmb{r}_-}-\mathrm{i}\pmb{a})^2\\
(a_x, a_y) = \frac{\theta}{\pi r^2}(-y, x) \tag{7.1.1}\\
\text{边界条件：}\ \ \phi(\pmb{r}_+,\pmb{r}_-) = \phi(\pmb{r}_+,-\pmb{r}_-)
\end{gather*}

当 $\theta=0$ 且 $\pmb{a}=0$ 时，$H$ 显然描述两个玻色子。对于 $\theta=\pi$，我们可以做一个规范变换
\begin{equation*}
\phi(\pmb{r}_+,\pmb{r}_-) \rightarrow \mathrm{e}^{\mathrm{i}f(\pmb{r}_-)}\phi(\pmb{r}_+,\pmb{r}_-), \qquad f(\pmb{r}) \equiv \arctan\left(\frac{x}{y}\right)
\end{equation*}

%FIG{7.6}: {二维谐振子势阱中两个任意子的能级 $E$ 与简并度 $D$（相对运动的）。}

由于 $\nabla f(\pmb{r}) = \frac{1}{r^2}(-y, x)$，我们发现 $\mathrm{e}^{-\mathrm{i}f}(\partial_{\pmb{r}_-}-\mathrm{i}\pmb{a})^2\mathrm{e}^{\mathrm{i}f} = (\partial_{\pmb{r}_-})^2$，于是哈密顿量变为
\begin{equation*}
\begin{aligned}
H &= -\frac{1}{2m}(\partial_{\pmb{r}_+})^2 - \frac{1}{2}\frac{1}{m}(\partial_{\pmb{r}_-})^2\\
\phi(\pmb{r}_+,\pmb{r}_-) &= -\phi(\pmb{r}_+,-\pmb{r}_-)\\
\text{或}\ \ \phi(\pmb{r}_1,\pmb{r}_2) &= -\phi(\pmb{r}_2,\pmb{r}_1)
\end{aligned}
\end{equation*}
（译注：原书第二行右端误印为 $-\phi(\pmb{r}_+,\pmb{r}_-)$，此处按文意改正。）
它描述两个自由费米子。一般地说，即使是 $N$ 个费米子，我们也可以自洽地把 $\pmb{a}$ 规范掉，并通过“边界条件”（即反对称条件）来描述费米统计。然而，对于 $N$ 任意子系统，我们无法把 $\pmb{a}$ 规范掉、再用“边界条件”来表示分数统计。

我们可以求解二维谐振子势阱中的两个任意子系统，得到如图 7.6 所示的能谱。可以清楚地看到，能谱如何从玻色情形连续地变到费米情形。

一个有趣的问题是：三维中存在任意子吗？三维中两个全同自由粒子的位形空间为 $V_+^{(3D)}\otimes V_-^{(3D)}$，其中
\begin{equation*}
V_-^{(3D)} = \left\{\pmb{r}_-\big|_{\pmb{r}_-\neq 0,\ \pmb{r}_-\sim-\pmb{r}_-}\right\}
\end{equation*}
这里 $V_-^{(3D)}$ 不是单连通的，整体相位存在。然而，$V_-^{(3D)}$ 中的不可缩回路由 $Z_2$ 标记。

%FIG{7.7}: {两个交换回路 $C_1$ 与 $C_{-1}$ 在三维中可以通过绕连接两粒子的轴转动而连续地相互变换。}

%FIG{7.8}: {把一个通量--电荷束缚态绕另一个通量--电荷束缚态移动一周所诱导的相位。}

这是因为沿一个方向的交换（由回路 $C_1$ 描述）与沿相反方向的交换（由回路 $C_{-1}$ 描述）属于同一类，因为两条回路 $C_{\pm 1}$ 可以连续地相互变形（见图 7.7）。如果我们给回路 $C_1$ 指派整体相位 $\mathrm{e}^{\mathrm{i}\theta}$，那么回路 $C_{-1}$ 将具有整体相位 $\mathrm{e}^{-\mathrm{i}\theta}$。由于两条回路 $C_{\pm 1}$ 仅相差一个局域相位为零的可缩回路，我们还有 $\mathrm{e}^{\mathrm{i}\theta}=\mathrm{e}^{-\mathrm{i}\theta}$。因此，$\theta$ 只能取如下两个可能的值：
\begin{equation*}
\theta = \left\{\begin{array}{cl}
0, & \text{玻色子}\\
\pi, & \text{费米子}
\end{array}\right.
\end{equation*}

以上的讨论表明，任意子是量子力学中的一种数学可能性。下一个问题是：任意子怎样才能出现在一些原本不含任意子的物理模型中？事实上，任意子可以作为玻色子/费米子系统中的激发出现。下面我们将考虑一个任意子作为电荷与磁通的束缚态而出现的模型。考虑一个由电荷为 1 的玻色子 $\phi$ 组成的系统。假设玻色子形成 $q$ 玻色子束缚态 $\Phi_q=(\phi)^q$，且 $\Phi_q$ 发生玻色凝聚。有效拉格朗日量为
\begin{align*}
\mathcal{L}_{\text{eff}} ={}& \frac{1}{2}\left|(\partial_\mu-\mathrm{i}q a_\mu)\Phi_q\right|^2 - \frac{1}{2}a\left|\Phi_q\right|^2 - \lambda\left|\Phi_q\right|^4 + \frac{1}{2g}\left(f_{\mu\nu}\right)^2 \tag{7.1.2}\\
&\qquad\qquad \langle\Phi_q\rangle \neq 0
\end{align*}
其中 $a_\mu$ 是与电荷玻色子耦合的 $U(1)$ 规范场。在凝聚态中，涡旋携带磁通 $\frac{2\pi}{q}$（见问题 7.1.2）。考虑涡旋 $\frac{2\pi}{q}$ 与原来的玻色子 $\phi$ 的束缚态。把一个束缚态绕另一个束缚态半周，我们得到一个相位
\begin{equation*}
\theta = \frac{1}{2}\frac{2\pi}{q} + \frac{1}{2}\frac{2\pi}{q} = \frac{2\pi}{q}
\end{equation*}
第一项来自电荷绕磁通管一周，第二项来自磁通管绕电荷一周（见图 7.8）。我们看到，$q=1$ 时束缚态是玻色子，$q=2$ 时是费米子，$q>2$ 时是任意子。

对于 BCS 超导体，$q=2$。但 $\phi$ 是费米子，因此束缚态是玻色子。目前所知唯一支持任意子激发的凝聚态系统是 FQH 系统。

\subsection*{问题 7.1.2}

在极坐标 $(r,\phi)$ 中，式 (7.1.2) 中的涡旋由
\begin{equation*}
\Phi_q = f(r)\mathrm{e}^{\mathrm{i}\phi}
\end{equation*}
描述，其中 $f(\infty)=\langle\Phi_q\rangle$。
\begin{enumerate}
\item[(a)] 证明：若 $a_\mu=0$，则涡旋的能量按 $\ln(L)$ 发散，其中 $L$ 是系统的尺寸。
\item[(b)] 求出使涡旋具有有限能量的 $\pmb{a}(r,\phi)$。
\item[(c)] 证明上述 $\pmb{a}(r,\phi)$ 的总磁通为 $2\pi/q$。
\end{enumerate}

\section{量子霍尔效应}

\subsection{整数量子霍尔效应}

\begin{keypoints}
\item 整数量子霍尔（IQH）态由填满的朗道能级描述。
\item $\nu=1$ IQH 态的多体波函数。
\item IQH 波函数的密度分布。
\end{keypoints}

让我们先讨论经典霍尔效应。考虑一个电荷 $q_e=-e$ 的电子二维气体，它在面内电场和垂直于该面的磁场的作用下运动。为了维持静态电流，电场力必须与磁场的洛伦兹力相平衡，即
\begin{equation*}
q_e\pmb{E} = q_e\pmb{v}\times\pmb{B}
\end{equation*}
在本章中，我们将选取使光速 $c=1$ 的单位。由于电流 $\pmb{j}$ 由 $\pmb{v}n$ 给出（$n$ 为密度），$\pmb{E}$ 与 $\pmb{j}$ 通过下式相联系：
\begin{equation*}
\pmb{E}\perp\pmb{j}, \qquad |\pmb{E}| = |\pmb{j}|\frac{B}{nq_ec} = |\pmb{j}|\frac{h}{q_e^2}\frac{B/(hc/q_e)}{n} = |\pmb{j}|\left(\frac{h}{q_e^2}\right)\frac{1}{\nu}
\end{equation*}
其中
\begin{equation*}
\nu \equiv \frac{nhc}{q_eB} = \frac{\text{粒子数}}{\text{磁通量子数}}
\end{equation*}
%FIG{7.9}: {霍尔电阻 $\rho_{xy}$ 随 $B$ 的变化。}
为填充因子。霍尔电阻 $\rho_{xy}=\left(\frac{h}{q_e^2}\right)\frac{1}{\nu}$ 可以在实验中直接测量。按照上述经典理论，若 $n$ 固定，则 $\rho_{xy}\propto B$。

实验上确实发现，对于二维电子气体，弱场下 $\rho_{xy}\propto B$（见图 7.9）。20 世纪 80 年代初，物理学家把二维电子气体置于强磁场（$\sim$10 特斯拉）中并冷却到极低的温度（$\sim$1 K），发现对于强 $B$，$\rho_{xy}$ 呈现平台结构，如图 7.9 所示（von Klitzing et al., 1980; Tsui et al., 1982）。这一现象称为 QH 效应。物理学家很快发现了 QH 态的许多惊人性质。最令人惊讶的是，平台恰好出现在与下式对应的那些 $\rho_{xy}$ 处：
\begin{equation*}
\nu = \underbrace{1, 2, 3, \cdots}_{\text{IQH}},\ \underbrace{\frac{1}{3}, \frac{2}{3}, \frac{2}{5}, \frac{3}{7}, \cdots}_{\text{FQH}}
\end{equation*}
$\rho_{xy}$ 的值如此精确和稳定，以致它已成为电阻的标准。

平台背后的物理稍微复杂一些，因为它涉及杂质。简而言之，如果二维电子气体在这些 $\nu$ 值处形成不可压缩液体\footnote{即所有带电激发都具有有限的能隙。}，我们就能理解这些填充因子处的平台。下面我们将看到，由于朗道能级结构，电子在 $\nu=1,2,3,\cdots$ 处确实形成不可压缩态。这使我们能够理解整数 $\nu$ 处的平台以及整数量子霍尔效应。

让我们忽略库仑相互作用，考虑强磁场 $B$ 中的二维自由电子气体。我们将忽略电子自旋，或者假设电子自旋已被极化。均匀磁场中的单个电子由下式描述：
\begin{gather*}
H = -\frac{1}{2m}(\partial_i-\mathrm{i}q_eA_i)^2, \qquad \hbar = c = 1 \tag{7.2.1}\\
(A_x, A_y) = \frac{B}{2}(-y, x), \qquad \text{（对称规范）}\\
B = \partial_xA_y - \partial_yA_x
\end{gather*}

上述哈密顿量有一个有趣的对称性质。尽管磁场是均匀的，它在直接平移 $T_{\pmb{a}}=\mathrm{e}^{-\pmb{a}\cdot\pmb{\partial}_{\pmb{x}}}$ 下却并非不变，即 $T_{\pmb{a}}HT_{\pmb{a}}^{\dagger}\neq H$。然而，在平移 $T_{\pmb{a}}$ 之后接着做规范变换
\begin{equation*}
G_{\pmb{a}} = \mathrm{e}^{\mathrm{i}\frac{1}{2}Bq_e(xa^y-ya^x)} \tag{7.2.2}
\end{equation*}
它就是不变的；或者更精确地说，
\begin{equation*}
G_{\pmb{a}}T_{\pmb{a}}H(G_{\pmb{a}}T_{\pmb{a}})^{\dagger} = H
\end{equation*}
组合变换 $G_{\pmb{a}}T_{\pmb{a}}$ 是 $H$ 的一个对称性，称为磁平移。不同方向的磁平移互不对易：
\begin{equation*}
(G_{\pmb{a}}T_{\pmb{a}})(G_{\pmb{b}}T_{\pmb{b}}) = \mathrm{e}^{-\mathrm{i}q_eB(a^xb^y-a^yb^x)}(G_{\pmb{b}}T_{\pmb{b}})(G_{\pmb{a}}T_{\pmb{a}}) \tag{7.2.3}
\end{equation*}
因此，尽管哈密顿量在 $x$、$y$ 两个方向都有磁平移对称性，$H$ 的本征态却不能用动量矢量 $(k_x,k_y)$ 来标记。注意，相位 $q_eB(a^xb^y-a^yb^x)$ 恰好是把电荷 $q_e$ 沿由 $\pmb{a}$、$\pmb{b}$ 张成的平行四边形 $P_{\pmb{a}\pmb{b}}$ 绕一周所得到的相位，即 $\mathrm{e}^{\mathrm{i}q_e\oint_{P_{\pmb{ab}}}\mathrm{d}\pmb{x}\cdot\pmb{A}}$。

为了求出式 (7.2.1) 的能级，我们引入复坐标 $z=x+\mathrm{i}y$，于是
\begin{equation*}
\left\{\begin{array}{ll}
\partial_z = \frac{1}{2}(\partial_x-\mathrm{i}\partial_y), & x = \frac{1}{2}(z+\bar{z}),\\
\partial_{\bar{z}} = \frac{1}{2}(\partial_x+\mathrm{i}\partial_y), & y = \frac{1}{2\mathrm{i}}(z-\bar{z}).
\end{array}\right.
\end{equation*}
用 $z$ 与 $\bar{z}$ 表示，我们有
\begin{equation*}
H = -\frac{1}{m}(D_zD_{\bar{z}} + D_{\bar{z}}D_z)
\end{equation*}
其中
\begin{equation*}
\left\{\begin{array}{ll}
D_z = \partial_z-\mathrm{i}q_eA_z, & A_z = \frac{1}{2}(A_x-\mathrm{i}A_y) = \frac{B}{4\mathrm{i}}\bar{z}\\
D_{\bar{z}} = \partial_{\bar{z}}-\mathrm{i}q_eA_{\bar{z}}, & A_{\bar{z}} = \frac{1}{2}(A_x+\mathrm{i}A_y) = -\frac{B}{4\mathrm{i}}z
\end{array}\right.
\end{equation*}
哈密顿量 $H$ 可以通过一个非幺正变换 $\mathrm{e}^{+|z|^2/4l_B^2}$ 来简化，其中 $l_B$ 是磁长度，定义为
\begin{equation*}
l_B^2 = \left|\frac{c\hbar}{q_eB}\right|
\end{equation*}
（注意 $2\pi l_B^2B = \frac{hc}{q_e} = \Phi_0$ 给出一个单位磁通量子。）我们发现
\begin{equation*}
H = \mathrm{e}^{-|z|^2/4l_B^2}\left[-\frac{1}{m}\left(\tilde{D}_z\tilde{D}_{\bar{z}} + \tilde{D}_{\bar{z}}\tilde{D}_z\right)\right]\mathrm{e}^{+|z|^2/4l_B^2}
\end{equation*}
其中
\begin{align*}
\tilde{D}_z &= \mathrm{e}^{+|z|^2/4l_B^2}D_z\mathrm{e}^{-|z|^2/4l_B^2} = \partial_z - \frac{q_eB}{4c}\bar{z} - \frac{1}{4l_B^2}\bar{z} = \partial_z - \frac{1}{2l_B^2}\bar{z}\\
\tilde{D}_{\bar{z}} &= \partial_{\bar{z}} + \frac{q_eB}{4c}z - \frac{1}{4l_B^2}z = \partial_{\bar{z}}
\end{align*}
这里我们假定了 $q_eB>0$。我们看到，非幺正变换 $\mathrm{e}^{+|z|^2/4l_B^2}$ 可以看作一个使 $A_{\bar{z}}=0$ 的“规范变换”。我们可以进一步把 $H$ 简化为
\begin{equation*}
H = \mathrm{e}^{-|z|^2/4l_B^2}\left[-\frac{2}{m}\left(\partial_z-\frac{1}{2l_B^2}\bar{z}\right)\partial_{\bar{z}}\right]\mathrm{e}^{+|z|^2/4l_B^2} + \frac{1}{2}\hbar\omega_c
\end{equation*}
其中 $\omega_c=q_eB/mc$ 是回旋频率。

现在我们可以容易地写出均匀磁场中电子的本征态：
\begin{equation*}
\begin{array}{ll}
\Psi_0 = \mathrm{e}^{-|z|^2/4l_B^2}f(z) & E = \frac{1}{2}\hbar\omega_c\\[5pt]
\Psi_1 = \mathrm{e}^{-|z|^2/4l_B^2}\left(\partial_z-\frac{1}{2l_B^2}\bar{z}\right)f(z) & E = \left(\frac{1}{2}+1\right)\hbar\omega_c\\[5pt]
\Psi_n = \mathrm{e}^{-|z|^2/4l_B^2}\left(\tilde{D}_z\right)^n f(z) & E = \left(\frac{1}{2}+n\right)\hbar\omega_c
\end{array}
\end{equation*}
这给出了朗道能级结构。

在上面的讨论中我们假定了 $q_eB>0$。第零朗道能级中的电子波函数是解析函数 $f(z)$ 乘以 $\mathrm{e}^{-|z|^2/4l_B^2}$。若 $q_eB<0$，则非幺正变换 $\mathrm{e}^{+|z|^2/4l_B^2}$ 将把 $A^z$ 变为零，第零朗道能级中的电子波函数将是反解析函数 $f(z^*)$ 乘以 $\mathrm{e}^{-|z|^2/4l_B^2}$。

第零朗道能级中的波函数 $\Psi_0$ 可以用基态
\begin{equation*}
\psi_m = \sqrt{N_m}\,z^m\,\mathrm{e}^{-|z|^2/4l_B^2}
\end{equation*}
展开，它们携带角动量 $m$。波函数 $\psi_m$ 呈环状。$|\psi_m| = \mathrm{e}^{m\ln|z|-|z|^2/4l_B^2}$ 的极大值位置在
\begin{equation*}
|z| = r_m = \sqrt{2m}\,l_B
\end{equation*}
我们发现第 $m$ 个态的环所围的面积为 $\pi r_m^2 = 2\pi l_B^2 m$，其中含有 $m$ 个磁通量子（见图 7.10）。

%FIG{7.10}: {第零朗道能级中的圆形轨道。}

%FIG{7.11}: {$\nu=1$ 液滴的密度分布，其中前 $m$ 个能级（由粗线表示）已被填满。}

因此，每一个磁通量子对应一个态，而第零朗道能级中的态数等于磁通量子数。于是，当 $\nu=1$ 时，第零朗道能级中的每个态都被一个电子占据。这给出有限的激发能隙 $\Delta=\hbar\omega_c$，并解释了实验中观察到的 $\nu=1$ 平台。不可压缩性来源于泡利不相容原理（费米统计）。

$\nu=1$ 态的多体波函数具有如下形式：
\begin{equation*}
\begin{aligned}
\Psi(z_1,\cdots z_N) &= \mathcal{A}\left[(z_1)^0(z_2)^1(z_3)^2\cdots\right]\mathrm{e}^{-\sum|z_i|^2/4l_B^2}\\
&= \det\left|\begin{array}{ccc}
1 & 1 & \ldots\\
z_1 & z_2 & \ldots\\
(z_1)^2 & (z_2)^2 & \ldots\\
\vdots & \vdots & \vdots
\end{array}\right|\mathrm{e}^{-\sum|z_i|^2/4l_B^2} = \prod_{i<j}(z_i-z_j)\,\mathrm{e}^{-\sum|z_i|^2/4l_B^2}
\end{aligned}
\end{equation*}
其中 $\mathcal{A}$ 是反对称化算符。上述波函数描述一个具有均匀密度的圆形液滴（图 7.11），因为每个电子占据相同的面积 $\pi r_{m+1}^2-\pi r_m^2 = 2\pi l_B^2$。这个性质很重要。一个一般的 $N$ 变量波函数在 $N\to\infty$ 极限下未必具有均匀的密度；那样的波函数不具有合理的热力学极限。

%FIG{7.12}: {$\nu=1/3$ 液滴的密度分布。}

\subsection*{问题 7.2.1}

求电子在势场 $V=\frac{1}{2}m\omega_0^2r^2$ 与均匀磁场 $B$ 中的能级。

\subsection*{问题 7.2.2}

\begin{enumerate}
\item[(a)] 证明：若规范变换 $G_{\pmb{a}}$ 由式 (7.2.2) 给出，则磁平移 $G_{\pmb{a}}T_{\pmb{a}}$ 是均匀磁场中哈密顿量 (7.2.1) 的一个对称性。
\item[(b)] 证明磁平移满足代数关系 (7.2.3)。
\end{enumerate}

\subsection{分数量子霍尔效应}

\begin{keypoints}
\item FQH 态的 Laughlin 波函数。
\item 等离子体类比与 Laughlin 态的密度分布。
\end{keypoints}

当 $0<\nu<1$ 时，第零朗道能级只是部分填充。对自由电子而言，这给出一个巨大的基态简并度 ${\sim}\,\frac{N_\phi!}{N!(N_\phi-N)!}$（其中 $N_\phi$ 是第零朗道能级中的态数，它同时也是磁通量子数；$N$ 是电子数）。因此，实验上观察到的 $\nu=1/3$ 态必定来自相互作用。这里有两个问题。第一，相互作用怎样才能产生能隙（不可压缩性）？第二，为什么 $\nu=1/3$ 是特别的？

Laughlin 用如下尝试波函数回答了上述两个问题：
\begin{equation*}
\Psi_3 = \prod_{i<j}(z_i-z_j)^3\,\mathrm{e}^{-\sum|z_i|^2/4l_B^2}
\end{equation*}
波函数 $\Psi_3$ 之所以好，有如下几条理由：
\begin{enumerate}
\item[a)] $\Psi_3$ 对每一对电子都有一个三阶零点 $(z_i-z_j)^3$。这对排斥型相互作用有利。
\item[b)] $\Psi_3$ 描述一个密度均匀的圆形液滴（图 7.12），其填充为 $\nu=1/3$。
\item[c)] $\Psi_3$ 是一个不可压缩态。
\end{enumerate}

由于 $|\Psi_3|$ 具有旋转不变性，显然 $|\Psi_3|$ 描述一个圆形液滴。如果我们再假设液滴还具有均匀的密度，那么就很容易
% ——本块到此为止：句子未完，接下一块 ch7_b（书页 288 / p303 顶部 “understand why it has a filling fraction ν=1/3.”），
% 由 ch7_b 的 A-TAIL 补出“理解它为何具有填充因子 ν=1/3。”
